\documentclass[12pt]{article}
\usepackage{amsmath}
\usepackage{natbib}

\title{Dual-Class Share Structures and the Cost of Equity Capital:
Evidence from Staggered Sunset Provisions\thanks{We thank seminar participants
for helpful comments.}}
\author{Jane Q. Researcher \and John A. Coauthor}
\date{\today}

\begin{document}
\maketitle

\begin{abstract}
We study how mandatory sunset provisions in dual-class share structures affect
firms' cost of equity capital. Using a panel of U.S. IPOs between 2005 and 2022
and exploiting staggered exchange-level rule changes that imposed time-based
sunsets, we estimate a difference-in-differences model. Firms subject to a
binding sunset experience a decline in their implied cost of equity of roughly
70 basis points relative to perpetual dual-class peers. The effect is
concentrated among founder-controlled firms with weak external monitoring.
Our findings speak to the debate over the efficiency of disproportionate
voting rights and the design of investor protection.
\end{abstract}

\section{Introduction}
The optimal allocation of voting rights between insiders and outside
shareholders is a central question in corporate governance \citep{grossman1988}.
Dual-class structures, which grant founders super-voting shares, have grown
rapidly among technology IPOs.\footnote{See \citet{bebchuk2017} for a discussion
of the rise of perpetual dual-class firms and the agency costs they create over
long horizons.} Proponents argue that they insulate visionary founders from
short-term market pressure, whereas critics emphasize the entrenchment and
private-benefit extraction they enable.

This paper asks whether \emph{sunset provisions} — clauses that automatically
convert super-voting shares to one-share-one-vote after a fixed period —
mitigate the valuation discount associated with entrenched control.

\section{Institutional Background and Hypotheses}
We describe the staggered adoption of time-based sunset requirements by major
U.S. exchanges and develop hypotheses linking the expected expiration of
control to the firm's cost of equity. The market-monitoring hypothesis predicts
that anticipated sunsets reduce expected agency costs.

\section{Data and Identification}
Our sample combines IPO prospectuses, hand-collected charter provisions, and
analyst forecasts used to compute the implied cost of equity following standard
accounting-based estimators. Identification rests on the staggered timing of
exchange rule changes, which we argue is plausibly exogenous to any single
firm's financing conditions.

\section{Results}
The baseline difference-in-differences estimate indicates a statistically
significant reduction in the implied cost of equity for treated firms. The
effect strengthens with proxies for weak monitoring and is robust to
alternative cost-of-equity measures and to controlling for time-varying
industry shocks.

\section{Conclusion}
Sunset provisions appear to lower the cost of equity for dual-class firms,
particularly where outside monitoring is weak. The evidence informs ongoing
policy debates over the regulation of disproportionate voting rights.

\bibliographystyle{plainnat}
\bibliography{references}

\end{document}
